\documentclass[10pt,a4paper]{amsart}
\usepackage[margin=19mm]{geometry}
\usepackage{amsmath,amssymb,mathtools}
\usepackage[scr=rsfs,cal=euler]{mathalfa}
\usepackage{xfrac}
\usepackage[unicode,hidelinks,bookmarks=false]{hyperref}
\newcommand{\ie}{i.e.\ }
\newcommand{\cf}{cf.\ }
\newcommand{\resp}{respectively\ }
\newcommand{\Subsection}{\subsection}
\providecommand{\nobreakdash}{\nobreak\hbox{-}}
\setlength{\emergencystretch}{2em}
\multlinegap0pt
\usepackage{bm}
\usepackage{bbm}
\usepackage{dsfont}
\usepackage{xspace}
\usepackage{cancel}
\usepackage{mathtools}
\usepackage{cleveref}
\usepackage{amsthm}
\makeatletter
\@ifundefined{theorem}{\newtheorem{theorem}{Theorem}}{}
\@ifundefined{claim}{\newtheorem{claim}[theorem]{Claim}}{}
\@ifundefined{lemma}{\newtheorem{lemma}[theorem]{Lemma}}{}
\makeatother
\def\PP{\mathbb{P}}
\def\T{\mathcal{T}}
\def\h{\hat}
\def\mcal{\mathcal}
\def\sufficient{matching}
\def\uhr{\upharpoonright}
\def\vwsamelong{\mathsf{OSub}^=}
\def\wsamelong{\mathsf{WSub}^{=}}
\def\wshorter{\mathsf{WSub}^{*}}
\newcommand{\NN}[0]{\mathbb{N}}
\newcommand{\OSub}{\mathsf{OSub}}
\newcommand{\RCA}[0]{\mathsf{RCA}}
\newcommand{\SubW}{\mathsf{Sub}}
\newcommand{\WSub}{\mathsf{WSub}}
\newcommand{\qvdash}{\operatorname{{?}{\vdash}}}
\title[The reverse mathematics of the Ordered Variable Word theorem]{The reverse mathematics of the Ordered Variable Word theorem}
\author{Lu Liu, Ludovic Patey}
\date{Source: arXiv:2606.12962v2 (CC BY 4.0)}
\begin{document}
\maketitle

\noindent\emph{Source and licence.} The reverse mathematics of the Ordered Variable Word theorem. Version arXiv:2606.12962v2 (CC BY 4.0), distributed under the Creative Commons Attribution 4.0 International licence, as recorded in the arXiv licence metadata for this exact version and shown on the abstract page https://arxiv.org/abs/2606.12962v2. Full rights notice: licenses/arxiv-2606.12962-CC-BY-4.0.txt.\par\smallskip

\noindent\textit{Editorial scope.} Counted unit: the Lemma whose source label is \texttt{None} in the article (source0.tex lines 1123--1133, source label \texttt{None}), delivered together with the article's own complete original proof of it (source0.tex lines 1134--1293, one \texttt{proof} environment of 160 lines). No mathematics is written, repaired or re-derived: statement and proof are the article's own text line for line. Required context retained uncounted: lem:graham-rothschild-tree (lines 565--569). Every definition, display and label of those lines is retained unchanged. Omitted from this package and not redistributed: 1--564, 570--1122, 1294--3353. Nothing inside a retained excerpt is omitted or altered: every retained body is delivered exactly as the article's lines are written, so no omission receipt of any kind accompanies this package. Citations: the retained excerpts carry no citation command at all, so no bibliography entry of the article is transcribed and no reference is redistributed. Reference adaptation: none was needed and none was made; every reference inside the delivered unit resolves inside it, so no unresolved reference or citation is produced. Printed slips kept literal: no printed word, symbol, hypothesis or number of the counted statement or of its proof was altered. Layout and class: the article's own class is \texttt{article} with packages including graphicx, amsthm, amsfonts, amsmath, amssymb, hyperref, cleveref, inputenc, url, multicol, enumitem, tikz, xcolor; the wrapper is the lane's pinned \texttt{amsart} editorial wrapper at 10pt on A4, which restates the theorem-like environments the retained text needs and adds the article's own macro definitions that the retained text needs (\texttt{\textbackslash NN}, \texttt{\textbackslash OSub}, \texttt{\textbackslash PP}, \texttt{\textbackslash RCA}, \texttt{\textbackslash SubW}, \texttt{\textbackslash T}, \texttt{\textbackslash WSub}, \texttt{\textbackslash h}, \texttt{\textbackslash mcal}, \texttt{\textbackslash qvdash}, \texttt{\textbackslash sufficient}, \texttt{\textbackslash uhr}, \texttt{\textbackslash vwsamelong}, \texttt{\textbackslash wsamelong}, \texttt{\textbackslash wshorter}), with extra wrapper packages (bm, bbm, dsfont, xspace, cancel, mathtools, cleveref). The delivered numbering is the unit's own. Assets: no retained span contains a figure environment, a graphics-inclusion command, a PDF-page inclusion, a table or a TikZ picture; no image file is copied into this package, the package declares no asset dependency and no image file is redistributed with it. Licence: version arXiv:2606.12962v2 of this article is distributed under the Creative Commons Attribution 4.0 International licence, as recorded in the arXiv licence metadata for this exact version and shown on the abstract page https://arxiv.org/abs/2606.12962v2. A full rights notice is delivered at licenses/arxiv-2606.12962-CC-BY-4.0.txt. Characters: every retained excerpt is pure ASCII, so the delivered engine, XeLaTeX with its default Latin Modern text font, reports no missing character.
\begin{lemma}[$\RCA_0$]\label[lemma]{lem:graham-rothschild-tree}
Fix~$A$ and $\ell \geq 1$.
Let $f : A^{<\omega} \to \ell$ be a coloring.
For every~$u \in A^{<\omega}$ and $n \in \NN$, there is some ordered $n$-variable word~$w$ at level~$n$ in~$\T_{A,\ell}$ such that $\WSub^{=}(w \cdot u)$ is $f$-monochromatic.
\end{lemma}

\begin{lemma}[$\Sigma_1^0$-Extension]\label[lemma]{caOVW-lemextSigma1}
Let $(I,W)$ be a $\PP$-condition and $J\subseteq I$.
Suppose for every $v\in J$,
$(I,W)\qvdash_v \varphi$.
Then there exists a $\PP$-condition $(\h I,\h W)\leq (I,W)$
such that for every $\h v\in \h I$ with $\h v$ extending
some $v\in J$, we have
$(\h v,\h W)\Vdash\varphi$.
Moreover, we can choose $(\h I,\h W)$
so that $\h W = W\uhr_{t(\h I)-t(I)}^\infty$  and the index of  $\h I $ can be obtained $W'$-computably.
\end{lemma}

\begin{proof}
 The main difficulty is to preserve
 the $f$-matching property while extending stems in $I$.
 This is possible since each $v\in J$ is extended
 to sufficiently many $\h v$. Indeed, we concatenate each such $v$ with some
 ordered variable word generated by nodes in some level of $\mcal T_{A, \ell \times I}$
 for \emph{every} node in that level.
 For the remaining stems $v\in I\setminus J$, we extend $v$ to some
 $v\cdot w$ where $w$ ranges over a set of words.
 The concrete proof is as follows.

Unfolding the definition of
$(I,W)\qvdash_v\varphi(G)$, there is a level $n \in \NN$
such that
for every $v\in J$,
every $u$ in level $n$ of  $\mcal T_{A, \ell \times I}$,
\begin{align}
\label{caOVW-eq16}
&\text{there exists some $\h u\in \vwsamelong(u) $
}\\ \nonumber
&\text{such that
 % }\\ \nonumber
%&\text{ %$ v\cdot\h u $ is \sectionallymonochromatic\ for $f$ and
$\varphi(v\cdot W[\h u])$ holds.}
\end{align}

In this proof,
let $witness(v)$ denote the set of all  $\h u$ satisfying (\ref{caOVW-eq16}).
Recall that by construction of the Graham-Rothschild tree, all the nodes at level $n$ in $\mcal T_{A, \ell \times I}$
are of length $t_n$.
Let
\begin{align}\label{caOVW-eq17}
\h I & := \left\{ v\cdot W[\hat u] : \begin{array}{l}
    \text{ either }v\in J\text{ and } \h u\in witness(v)\\
  \text{ or } v\in I\setminus J\text{ and }\hat u\in A^{t_n}
  \end{array} \right\} \\
\h W & :=  W \uhr_{t(\h I)-t(I)}^\omega. \nonumber
\end{align}
We show that $(\h I,\h W)$ is the desired extension.
It is obviously a $\PP$-precondition.
By definition of $\h I$ (see also (\ref{caOVW-eq16})),
for every $\h v\in \h I$ with $\h v$ extending
some member of $J$, $\varphi(\h v)
$ holds (so $(\h v,\h W)\Vdash\varphi(G)$).
Thus, it remains to show that
\begin{claim}\label[claim]{caOVW-claim4}
$(\h I,\h W)$ is $f$-\sufficient.
\end{claim}
\begin{proof}

Let $$w\in \wshorter(W\uhr_{t(\h I)-t(I)}^\infty).$$
It suffices to find $\h v\in \h I$ such that
\begin{align}\label{caOVW-eq20}
\wsamelong(\h v\cdot w)\text{ is monochromatic for $f$.}
\end{align}
Since $(I, W)$ is $f$-matching, let $\bar f : A^{<\omega} \to \ell \times I$ be defined as follows:
\begin{align}\label{newcaOVW-eq1}
\bar f: &\ A^{<\omega}\ni w \mapsto \langle c, v\rangle
\text{ such that }\\ \nonumber
&\text{$\wsamelong(v\cdot W[w])$ is $f$-monochromatic for color~$c$}.
\end{align}
Let $\hat w \in A^{<\omega}$ be such that
$$w = W\uhr_{t(\h I)-t(I)}^\omega[\hat w].$$
By the fundamental property of the Graham-Rothschild tree $\T_{A,\ell \times I}$ (\Cref{lem:graham-rothschild-tree}) applied to the coloring~$\bar f$ and the word $\hat w$, there is some~$u$ at level~$n$ in $\T_{A,\ell \times I}$ such that
\begin{align}\label{newcaOVW-eq2}
\text{$\WSub^=( u \cdot \hat w)$ is $\bar f$-monochromatic for some color $\langle c, v\rangle \in \ell \times I$} 
\end{align}
By (\ref{newcaOVW-eq1}),
\begin{align}\label{newcaOVW-eq3}
\text{$\WSub^=(v \cdot W[u \cdot \hat w])$ is $f$-monochromatic for color~$c$} 
\end{align}
Note that
$$
W[u \cdot \hat w] = W[u] \cdot W\uhr_{t(\h I)-t(I)}^\omega[\hat w] = W[u] \cdot w
$$
So by (\ref{newcaOVW-eq3}), we have
\begin{align}\label{newcaOVW-eq4}
\text{$\WSub^=(v \cdot W[u] \cdot w)$ is $f$-monochromatic for color~$c$} 
\end{align}
Suppose first that $v \in I \setminus J$. In particular, letting $\hat u \in \SubW^=(u)$,
we have in particular $\hat u \in A^{t_n}$, so by (\ref{caOVW-eq17}), $v \cdot W[\hat u] \in \hat I$. Moreover, by (\ref{newcaOVW-eq4}), we have
$$
\text{$\WSub^=(v \cdot W[\hat u] \cdot w)$ is $f$-monochromatic for color~$c$.} 
$$
Suppose now that $v \in J$. By (\ref{caOVW-eq16}), there is some $\hat u \in witness(v)$ such that $\hat u \in \OSub^=(u)$. By (\ref{caOVW-eq17}), $v \cdot W[\hat u] \in \hat I$. Moreover, by (\ref{newcaOVW-eq4}), we have
$$
\text{$\WSub^=(v \cdot W[\hat u] \cdot w)$ is $f$-monochromatic for color~$c$.} 
$$
In both cases we are done. This completes the proof of \Cref{caOVW-claim4}.
% \ludovic{TODO}


% Note that $W\uhr t(\h I)-t(I)$
% contains $n$ many variable kinds;
% and for every $\sigma\in A^{t_n}$,
% $W[\sigma]\cdot w\in \wshorter(W)$.
% So by $f$-\sufficient\ of $(I,W)$, for every $\sigma\in A^{t_n}$,
% there exists $v\in I$ such that
% \begin{align}\label{caOVW-eq21}
% \text{$\wsamelong(v\cdot W[\sigma]\cdot w)$ is monochromatic for $f$.}
% \end{align}
% If,  for some $\sigma\in A^{t_n} $ and $v\in I\setminus J$,
% $$
% \text{$\wsamelong(v\cdot W[\sigma]\cdot w)$ is monochromatic for $f$,}
% $$
% then we are done  verifying
% (\ref{caOVW-eq20}) since $v\cdot W[\sigma]\in \h I$
% (see (\ref{caOVW-eq17})).
% Therefore assume otherwise.
% Now (\ref{caOVW-eq21}) becomes:
%  for every $\sigma\in A^{t_n}$,
% there exists $v\in   J $ (instead of $v\in I$) such that
% \begin{align}\nonumber
% \text{$\wsamelong(v\cdot W[\sigma]\cdot w)$ is monochromatic for $f$.}
% \end{align}
% In another word, $w$ give rise to the following coloring $A^{t_n}\rightarrow  \ell \times J$,
% \begin{align}\label{caOVW-eq112}
% g: &A^{t_n}\ni \sigma\mapsto   \langle \h\ell, v\rangle \in \ell \times J
% \text{ iff }\\ \nonumber
% &\text{$\wsamelong(v\cdot W[\sigma]\cdot w)$ is monochromatic for $f$
% for color $\h \ell$}.
% \end{align}
% By construction of \GRtree\ $\mcal T_{A, \ell \times I}$,
% there exists a $u$ in level $t$ of $\mcal T_{A, \ell \times I}$ such that
% $$
% \wsamelong(u)\text{ is monochromatic for $g$, say for color $ \langle \ell^*, v^*\rangle $.}
% $$
% Unfolding definition of $g$ (see (\ref{caOVW-eq112})),
% we have
% \begin{align}\label{caOVW-eq113}
% \text{$\wsamelong(v^*\cdot W[u]\cdot w)$ is monochromatic for $f$
% (for color $\ell^*$).}
% \end{align}

% By (\ref{caOVW-eq16})(\ref{caOVW-eq17}), let
% \begin{align}\label{caOVW-eq18}
% \h u\in \vwsamelong(u)
% \text{ be such that }
% v^*\cdot W[\h u]\in \h I.
% \end{align}
% By (\ref{caOVW-eq113})
% and since
% $
% \h u\in \vwsamelong(u),
% $
% we have
% \begin{align} \nonumber
% \text{$\wsamelong(v^*\cdot W[\h u]\cdot w)$ is monochromatic for $f$.}
% \end{align}
% Combine with
% $ v^*\cdot W[\h u]\in \h I$ (see (\ref{caOVW-eq18}))
%  we  verified (\ref{caOVW-eq20}).
%   Thus, we are done proving Claim \ref{caOVW-claim4}.



\end{proof}
The ``Moreover" part follows by analyzing  the effectiveness of above proof.
This completes the proof of \Cref{caOVW-lemextSigma1}.
\end{proof}

\end{document}
